\par\noindent\textbf{(1)} 第 \(n\) 格点的合力为 \(K(u_{n+1}+u_{n-1}-2u_n)-\kappa u_n\)。Newton方程 \(m\ddot u_n=K(u_{n+1}+u_{n-1}-2u_n)-\kappa u_n\)，能量 \(H=\sum_n[m\dot u_n^2/2+K(u_{n+1}-u_n)^2/2+\kappa u_n^2/2]\) 为正定。
\par\noindent\textbf{(2)} 设 \(u_n=Ae^{i(qna-\omega t)}\)，则 \(\omega^2(q)=\kappa/m+(4K/m)\sin^2(qa/2)\)。取正频支，群速 \(v_g=d\omega/dq=(Ka/(m\omega))\sin(qa)\)。当 \(q=0\) 因钉扎存在频隙 \(\omega_{\min}=\sqrt{\kappa/m}\)，不是声学零模。
\par\noindent\textbf{(3)} 对于 \(|qa|\ll1\)，\(\sin(qa/2)\approx qa/2\)，故 \(\omega^2\approx\omega_0^2+c^2q^2\)，其中 \(\omega_0^2=\kappa/m\)、\(c^2=Ka^2/m\)。连续位移 \(u(x,t)\) 满足 \(u_{tt}=c^2u_{xx}-\omega_0^2u\)，即一维Klein–Gordon方程。
\par\noindent\textbf{(4)} 在 \(q=\pm\pi/a\) 时 \(\omega_{\max}=\sqrt{(\kappa+4K)/m}\) 且 \(v_g=0\)，因此包络在带边不传播。若 \(\kappa\to0\)，低波数恢复 \(\omega\sim c|q|\) 的声学分支，整数格点整体平移不消耗弹性能。
